\chapter*{Summary}
\addcontentsline{toc}{chapter}{Summary}
\setheader{Summary}

This thesis presents the design, implementation, and evaluation of CryptoVolt, an AI-assisted algorithmic trading system that combines multi-source sentiment analysis with conventional technical indicators and gradient-boosted learning to produce risk-aware trading decisions on cryptocurrency markets. The central problem addressed by this work is that many automated systems rely almost exclusively on price history and therefore respond poorly when information shocks, rather than smooth technical patterns, dominate short-term dynamics. CryptoVolt addresses this gap by giving sentiment a structural role in the decision pipeline, including an execution-level veto when informational risk is elevated.

The system is organised as a FastAPI backend with a React (Vite) dashboard, optional Go-based desktop agent for local repository analytics, and persistent storage via SQLAlchemy with either SQLite or PostgreSQL. Machine learning centres on XGBoost classifiers trained from engineered features, with sequence models implemented in PyTorch where configured. Sentiment scoring uses VADER together with Reddit ingestion when API credentials are available. A hybrid decision endpoint fuses rule-based logic, model outputs, and veto reasoning; paper trading routes ensure experimentation without live capital.

The platform targets Binance-compatible perpetual-style symbols and operates in paper trading mode by default. Discord webhooks can deliver operator notifications when configured. Evaluation combines offline backtesting, classifier metrics on held-out data, paper-trading statistics, and analysis of veto events. Reported headline figures include improved Sharpe ratios and reduced maximum drawdown relative to baselines, together with a high fraction of vetoes that prevented hypothetical losses during stress windows. These results should be interpreted within the limitations of simulation assumptions, dataset scope, and the English-heavy sentiment sources used in this prototype.
